\documentclass{article}

\usepackage{amsmath}
\usepackage{amsfonts}
\usepackage{amssymb}
\usepackage{multicol}

\usepackage{mathenv}

\def\nbOne{{\mathchoice {\rm 1\mskip-4mu l} {\rm 1\mskip-4mu l}
{\rm 1\mskip-4.5mu l} {\rm 1\mskip-5mu l}}}

\usepackage{vmargin}
\setmarginsrb{2.5cm}{2.5cm}{2.5cm}{2.5cm}{0cm}{0cm}{0cm}{0cm}

\usepackage[utf8]{inputenc}

\usepackage[french]{babel}
\selectlanguage{french}

\usepackage{color}
\usepackage{hyperref}
\hypersetup{pdfborder={0 0 0}, colorlinks=true, urlcolor=blue, linkcolor = darkred}
\usepackage{graphicx}
\graphicspath{{img/}} 
\usepackage{listings}
\definecolor{colKeys}{rgb}{0.75,0,0}
\definecolor{colIdentifier}{rgb}{0,0,0}
\definecolor{colComments}{rgb}{0.75,0.75,0}
\definecolor{colString}{rgb}{0,0,0.7}

\usepackage{verbatim}
\usepackage{moreverb}

\lstset{
basicstyle=\ttfamily\small, %
identifierstyle=\color{colIdentifier}, %
keywordstyle=\color{colKeys}, %
stringstyle=\color{colString}, %
commentstyle=\color{colComments}, %
showspaces=false,
}
\lstset{language=java}

% Commandes personnelles %

\definecolor{darkred}{rgb}{0.85,0,0}
\definecolor{darkblue}{rgb}{0,0,0.7}
\definecolor{darkgreen}{rgb}{0,0.6,0}
\definecolor{darko}{rgb}{0.93,0.43,0}
\definecolor{maintitle}{rgb}{0.66,0,0.22}
\definecolor{title}{rgb}{0,0.5,0.5}
\newcommand{\maintitlecolor}[1]{\textcolor{maintitle}{#1}}
\newcommand{\titre}[1]{\textcolor{title}{#1}}
\newcommand{\tsect}[1]{\titre{\section{#1}}}
\newcommand{\tssect}[1]{\titre{\subsection{#1}}}
\newcommand{\tsssect}[1]{\titre{\subsubsection{#1}}}
\newcommand{\vect}[1]{\overrightarrow{#1}}
\newcommand{\dred}[1]{\textcolor{darkred}{\textbf{#1}}}
\newcommand{\dgre}[1]{\textcolor{darkgreen}{\textbf{#1}}}
\newcommand{\dblu}[1]{\textcolor{darkblue}{\textbf{#1}}}
\newcommand{\dora}[1]{\textcolor{darko}{\textbf{#1}}}
\newcommand{\gre}[1]{\textcolor{darkgreen}{#1}}
\newcommand{\blu}[1]{\textcolor{darkblue}{#1}}
\newcommand{\ora}[1]{\textcolor{darko}{#1}}
\newcommand{\rouge}[1]{\textcolor{darkred}{#1}}
\newcommand{\ceil}[1]{\left\lceil #1 \right\rceil}
\newcommand{\cdil}[1]{\left\lfloor #1 \right\rfloor}
\newcommand{\term}[1]{\textit{\textcolor{maintitle}{#1}}}
\newcommand{\image}[1]{\includegraphics{#1}}
\newcommand{\imageR}[2]{\includegraphics[width=#2px]{#1}}
\newcommand{\imageRT}[2]{\includegraphics[height=#2px]{#1}}
\newcommand{\img}[1]{\begin{center}\includegraphics[width=400px]{#1}\end{center}}
\newcommand{\imag}[1]{\begin{center}\includegraphics{#1}\end{center}}
\newcommand{\imgR}[2]{\begin{center}\includegraphics[width=#2px]{#1}\end{center}}
\newcommand{\imgRT}[2]{\begin{center}\includegraphics[height=#2px]{#1}\end{center}}
\newcommand{\point}[2]{\item \ora{\underline{#1}} : \textit{#2}}
\newcommand{\bfp}[2]{\item \textbf{#1} : \textit{#2}}
\newcommand{\sumparam}[3]{\sideset{}{_{#1}^{#2}}\sum{#3}}
\newcommand{\sumin}[3]{\sideset{}{_{i=#1}^{#2}}\sum{#3}}
\newcommand{\sumkn}[3]{\sideset{}{_{k=#1}^{#2}}\sum{#3}}
\newcommand{\intin}[3]{\sideset{}{_{#1}^{#2}}\int{#3}}
\newcommand{\stitre}[1]{\noindent\textbf{\underline{#1}} \\}
\newcommand{\R}{\mathbb{R}}
\newcommand{\Z}{\mathbb{Z}}
\newcommand{\N}{\mathbb{N}}
\newcommand{\ualpha}{\vect{u_\alpha}}
\newcommand{\valpha}{\vect{v_\alpha}}
\newcommand{\palpha}{\vect{\Psi_\alpha}}
\newcommand{\npcomp}{\term{$\mathcal{NP}$-complet}}
\newcommand{\npcompl}{\term{$\mathcal{NP}$-complet} }
\DeclareMathAlphabet{\mathpzc}{OT1}{pzc}{m}{it}

\begin{sffamily}

\title{$ $\\ $ $\\ $ $\\ $ $\\ $ $\\ $ $\\ $ $\\\begin{Huge}\maintitlecolor{Intelligence Artificielle}\end{Huge} \\ 
   $ $ \\ \begin{LARGE}\textit{Devoir 1}\end{LARGE}}
\author{\textit{Xavier Dubuc} \\ \blu{Xavier.DUBUC@student.umons.ac.be} \\ $ $\\ $ $\\ $ $\\ $ $\\ $ $\\ $ $\\ $ $\\ $ 
$ \\ 
\includegraphics{UMONS.jpg}}
%\date{}
\end{sffamily}

\begin{document}\begin{sffamily}

\maketitle

\newpage

\tableofcontents

\hbox{\raisebox{0.4em}{\vrule depth 0.4pt height 0.4pt width 10cm}}

\newpage

\section{Question 1}

\term{Parcourez la littérature portant sur l'\textbf{IA} pour savoir si les tâches suivantes peuvent être actuellement
résolues par ordinateur : }
\begin{itemize}
\item \textbf{jouer un bon niveau au tennis de table,} le robot \textbf{Topio} développé par \textbf{TOSY} (firme de robotique 
vietnamienne) permet de jouer au tennis de table de manière assez convaicante, on peut voir la version \textbf{2.0} en oeuvre ici 
: \href{http://www.youtube.com/watch?v=SE6tFX5QM\_U}{http://www.youtube.com/watch?v=SE6tFX5QM\_U} , on l'aperçoit même en train de 
battre son adversaire. De plus, on peut apprendre sur \href{http://en.wikipedia.org/wiki/TOPIO}{wikipédia} \textit{(entre autres)} 
qu'une version \textbf{3.0} a été développée et qu'il apprend constamment en jouant, ce qui laisse à penser que cette version joue
bien mieux que la version \textbf{2.0}. La réponse est donc \gre{OUI}.
\item \textbf{conduire à Victorville, en Californie,}
Je cite \href{http://www.cs.kuleuven.ac.be/~dannyd/CrashlessCars.pdf}{ce pdf} : \textit{«In 2007 a tricked-out Chevrolet Tahoe 
nicknamed “Boss” and several similar driverless vehicles successfully navigated through a realistic city streetscape in 
Victorville, Calif., one complete with other cars and even traffic jams. The autonomous cars and trucks were competing in the 
Defense Advanced Research Projects Agency’s Urban Challenge, arace designed to demonstrate that robot road vehicles can become 
practical. Soon afterward General Motors CEO G.Richard Wagoner, Jr., predicted his company will market autonomous vehicles within 
10 years. That prognostication may be a bit optimistic, but his statement surely points the way to real robotic cars in the not 
too distant future.»} $\Rightarrow$ La réponse est donc \gre{OUI}.
\item \textbf{jouer au Bridge \textit{(niveau compétition)},}
l'article le plus récent, 
\href{http://www.sciencedirect.com/science?_ob=ArticleURL&_udi=B6TYF-3TP12X4-3&_user=532109&_coverDate=04/30/1998&_rdoc=1&_fmt=high&_orig=search&_origin=search&_sort=d&_docanchor=&view=c&_searchStrId=1508536596&_rerunOrigin=scholar.google&_acct=C000026718&_version=1&_urlVersion=0&_userid=532109&md5=1fc89a599671794b51c6346da562e219&searchtype=a}
{\textit{«Search in games with incomplete information: a case study using Bridge card play»}}, écrit par
\textbf{Ian Frank} et \textbf{David Basin}, fait état d'une intelligence artificielle qu'ils ont développé et qui permet de jouer 
au \textbf{Bridge} avec la tactique \textit{«Best Defense»}. On en déduit que le niveau compétition n'est pas atteint par cette
intelligence artificielle, la réponse est donc \rouge{NON}.
\item \textbf{découvrir et prouver de nouveaux théorèmes mathématiques,}
malgré de longues recherches, le peu que j'ai pu trouver fut quelques articles 
\href{http://www.sciencedirect.com/science?_ob=ArticleURL&_udi=B6TYF-47YJHBM-7&_user=532109&_coverDate=08/31/1979&_alid=1508526304&_rdoc=1&_fmt=high&_orig=search&_origin=search&_zone=rslt_list_item&_cdi=5617&_sort=r&_st=13&_docanchor=&view=c&_ct=109&_acct=C000026718&_version=1&_urlVersion=0&_userid=532109&md5=90dc57b7279c2f524a73f59c197bb90d&searchtype=a}
{des années 70} spécifiant comment prouver/réfuter un large nombre de résolution linéaire \textit{(même avec la VPN je n'ai pas pu 
avoir accès au fichier PDF)}. Rien n'est trouvable quant à la découverte de théorèmes par contre, ce qui oriente ma réponse vers 
le \rouge{NON}.
\end{itemize}

\section{Question 2}

\newcommand{\valeur}[1]{\gre{\textit{#1}}}

\begin{itemize}
\item \term{Décrire une fonction pour un agent rationnel dans le cas où chaque mouvement coûte un point (pour le 
reste les hypothèses sont identiques à celles données au transparent 22 du chapitre 2, y compris le gain d'un point
par carré néttoyé). Le programme de l'agent correspondant exige-t-il de conserver un état interne (agent avec 
modèle)?}\\


$\Rightarrow$\textbf{ Première ébauche}\\

\begin{center}\fbox{
\begin{minipage}{0.7\textwidth}
\blu{function} \rouge{MODEL-VACUUM-AGENT} ([\valeur{location},\valeur{status}]) \blu{returns} an action \\
\indent $\qquad$\blu{persistent} : \valeur{action}, l'action la plus récente (initialement vide) \\
\indent $\qquad$\blu{if} \valeur{status} = \valeur{Sale} \blu{then} \valeur{action} $\leftarrow$ \valeur{Aspire} \\
\indent $\qquad$\blu{else} \\
\indent $\qquad\qquad$\blu{if} \valeur{location} = \valeur{A} \blu{then} \\
\indent $\qquad\qquad\qquad$\blu{if} \valeur{action} = \valeur{Gauche} \blu{then} \valeur{action} $\leftarrow$ 
\valeur{Stop} \\
\indent $\qquad\qquad\qquad$\blu{else} \valeur{action} $\leftarrow$ \valeur{Droite} \\
\indent $\qquad\qquad$\blu{else} \\
\indent $\qquad\qquad\qquad$\blu{if} \valeur{action} = \valeur{Droite} \blu{then} \valeur{action} $\leftarrow$ 
\valeur{Stop} \\
\indent $\qquad\qquad\qquad$\indent \blu{else} \valeur{action} $\leftarrow$ \valeur{Gauche} \\
\indent $\qquad$\blu{return} \valeur{action}
\end{minipage}
}
$ $\\ 
$ $\\ 
\end{center}

Le robot se comporte de la manière suivante :
\begin{itemize}
\item si le carré où il se trouve est sale, quoi qu'il arrive, il l'aspire ;
\item si ce carré n'est pas sale, il consulte la dernière action qu'il a effectué, s'il vient de se déplacer, c'est que
le précédent carré était propre, il peut donc s'arrêter, son travail est terminé ;
\item sinon, dans tous les autres cas, (c'est-à-dire que le carré est propre et il ne vient pas de se déplacer), il 
doit changer de carré. \\
\end{itemize}

\textbf{Avec cette fonction d'agent}, le score maximal atteignable est 0 (par exemple, départ en A, il aspire (+1) il va 
en B (-1), il aspire (+1) puis retourne en A (-1) et s'arrête) ; pour maximiser ce score il faudrait empêcher le retour 
en A à la fin ; autrement dit stocker l'état du "monde". On va donc stocker l'état de cette manière : 
(\valeur{statusA},\valeur{statusB}). Ceci nous permettra d'éviter le dernier retour en A car \valeur{state} contiendra 
(\valeur{Propre},\valeur{Propre}) (cette variable contient (\valeur{Sale},\valeur{Sale}) initialement). \\

Voici cette fonction :\\

\newcommand{\fesp}{\indent $\qquad$}
\begin{center}\fbox{
\begin{minipage}{0.8\textwidth}
\blu{function} \rouge{MODEL-VACUUM-AGENT} ([\valeur{location},\valeur{status}]) \blu{returns} an action \\
\indent $\qquad$\blu{persistent} : \valeur{action}, l'action la plus récente (initialement vide) \\
\indent $\qquad\qquad\ \qquad$ \valeur{state}, contenant un couple représentant l'état de l'environnement \\
\fesp \blu{if} \valeur{status} = \valeur{Sale} \blu{then} \valeur{action} $\leftarrow$ \valeur{Aspire}\\
\fesp \blu{else} \\
\fesp $\qquad$ \blu{if} \valeur{location} = \valeur{A} \blu{then} \\
\fesp $\qquad\qquad$ \valeur{state} = (\valeur{Propre},\valeur{state}[1])\\
\fesp $\qquad\qquad$ \blu{if} \valeur{state} = (\valeur{Propre},\valeur{Propre}) then \valeur{action} $\leftarrow$ \valeur{Stop} \\
\fesp $\qquad\qquad$ \blu{else} \valeur{action} $\leftarrow$ \valeur{Droite}\\
\fesp $\qquad$ \blu{else} \\
\fesp $\qquad\qquad$ \valeur{state} = (\valeur{state}[0],\valeur{Propre})\\
\fesp $\qquad\qquad$ \blu{if} \valeur{state} = (\valeur{Propre},\valeur{Propre}) then \valeur{action} $\leftarrow$ \valeur{Stop} \\
\fesp $\qquad\qquad$ \blu{else} \valeur{action} $\leftarrow$ \valeur{Gauche}\\
\fesp \blu{return} \valeur{action}
\end{minipage}
}
$ $\\ 
$ $\\ 
$ $\\ 
\end{center}

Dans ce cas-ci, le robot ne se déplacera que s'il n'est pas sûr que l'autre carré n'est pas propre. Ainsi, dans l'exemple employé 
précédemment :
\begin{itemize}
\item il commence en \textbf{A}, \valeur{state} = (\valeur{Sale},\valeur{Sale}) et \valeur{status} = \valeur{Sale}, donc il aspire 
(+1)
\item \valeur{state} = (\valeur{Sale},\valeur{Sale}) et \valeur{status} = \valeur{Propre}, il met \valeur{state} à jour et se 
déplace en \textbf{B} (-1)
\item il arrive en \textbf{B}, \valeur{state} = (\valeur{Propre},\valeur{Sale}) et \valeur{status} = \valeur{Sale}, donc il aspire
(+1)
\item \valeur{state} = (\valeur{Propre},\valeur{Sale}) et \valeur{status} = \valeur{Propre}, il met \valeur{state} à jour et 
s'arrête car \valeur{state} = (\valeur{Propre},\valeur{Propre})
\end{itemize}

$\Rightarrow$ Le robot obtiendra donc un score compris entre $-1$ (les 2 carrés étaient propres) et $1$ (les 2 carrés étaient 
sales).

\textbf{Comme l'explication précédente le pointe, le robot nécessite la conservation d'un état interne lui permettant de savoir
à tout instant la dernière action effectuée.} \\

\item \term{Discutez des conceptions possibles d'un agent pour les cas où des carrés propres peuvent devenir sales
(par exemple, pour chaque unité de temps, chaque carré propre à 10\% de chances de devenir sale) et où la géographie
de l'environnement est inconnue. Est-il judicieux pour l'agent d'apprendre de son expérience dans ces cas ? Si oui,
que devrait-il apprendre ? Si non, pourquoi pas ?} \\

$\Rightarrow$ Il est très judicieux d'apprendre de son environnement, ne fusse que pour la géographie, le fait de savoir comment se
découpe son environnement peut l'amener à optimiser sa performance. Si on prend le même exemple que ci-dessus, sans connaître la
géographie de l'endroit où il évolue, ses décisions ne peuvent être aussi précises, il devra scruter son environnement à l'aveugle.
Quant au fait que les carrés redeviennent propres à 10\% de chances, il peut, en apprenant ça se dire qu'il ne revisite un carré 
que si ça fait $n \geq 10$ unités de temps qu'il ne l'a pas visité, minimisant ainsi ses déplacements (s'il y a 3 carrés et qu'il 
revisite sans cesse les carrés pour vérifier qu'ils sont toujours propres, il fait énormément de déplacements inutiles, alors que 
s'il aspire les 3 à la suite toutes les 10 unités de temps, il ne fait pas de déplacements inutiles).
\end{itemize}

\section{Question 3}

\term{Considérons un simple thermostat qui allume un four lorsque la température est d'au moins 3 degrés en dessous 
du réglage, et qui l'éteint lorsque la température est d'au moins 3 degrés au dessus du réglage. Un thermostat est-il 
une instance d'un simple agent réflexe, d'un agent réflexe basé sur un modèle ou d'un agent basé sur des objectifs ? 
Justifiez.} \\

$\Rightarrow$ Il s'agit d'un simple agent réflexe, sa décision ne dépend que de la perception courante. En effet, à chaque 
unité de temps, il obtient une perception de la température du four, si elle est supérieur au réglage il coupe le four, si elle 
est inférieur il l'allume sinon il ne fait rien.

\section{Répartition du temps}

\begin{enumerate}
\item entre $1$ et $2$h, certains sujets sont difficiles à trouver et les articles sur le bridge et Topio m'ont assez bien 
captivés, j'ai donc passé pas mal de temps à les lire (ou à regarder des vidéos).
\item $1$h environ, en comptant la mise en page LaTeX.
\item $5$ minutes.
\end{enumerate}

\end{sffamily}\end{document}